#import "@local/bootlin:0.1.0": *

#import "@local/bootlin-yocto:0.1.0": *

#import "@local/bootlin-utils:0.1.0": *

#import "../../typst/local/themeBootlin.typ": *

#import "../../typst/local/common.typ": *

#show: bootlin-theme.with( aspect-ratio: "16-9", config-common(handout: "handout" in sys.inputs and sys.inputs.handout == "1", ))

#show raw.where(block: true): set block(fill: luma(240), inset: 1em, radius: 0.5em, width: 100\%)

#show raw.where(block: true): set text(size: 11pt)

#show raw.where(block: false): r => text(fill: color-link)[#r]

#show raw.where(lang: "c", block: true): set block(fill: luma(240), inset: 0.4em, radius: 0.5em, width: 95\%, breakable: true, above: 12pt, below: 12pt)

#show raw.where(lang: "c", block: true): set text(11pt)

#show raw.where(lang: "console", block: true):set block(fill: luma(240), inset: 0.4em, radius: 0.5em, width: 95\%, breakable: true, above: 6pt)

#show raw.where(lang:"console", block: true): set text(12pt)

 = Useful resources

=== Books
  
#table(columns: (50\%, 50\%), stroke: none,
[
    \small
    \begin{itemize}
    \item {\bf Mastering Embedded Linux Programming, 4th Edition}
      \footnote{\tiny
\url{https://www.amazon.com/dp/1803232595}}\\
      By Frank Vasquez, Chris Simmonds, Packt Publishing, May 2025\\
      An up-to-date resource covering most aspects of embedded Linux
      development.
    \item {\bf The Linux Programming Interface}
      \footnote{\tiny \url{https://man7.org/tlpi/}}\\
      Michael Kerrisk (maintainer of Linux manual pages), 2010, No Starch Press\\
      A gold mine about Linux system programming\\
    \end{itemize}
    \normalsize
    \column{0.15\textwidth}
    \includegraphics[height=0.25\textheight]{slides/sysdev-references/book-mastering-embedded-linux4.jpg}\\
    \vspace{0.5cm}
    \includegraphics[height=0.25\textheight]{common/linux-programming-interface.png}\\
  \\]


=== Web sites
  \begin{itemize}
  \item {\bf ELinux.org}, \url{https://elinux.org}, a Wiki entirely
    dedicated to embedded Linux. Lots of topics covered: real-time,
    filesystems, multimedia, tools, hardware platforms,
    etc. Interesting to explore to discover new things.
  \item {\bf LWN}, \url{https://lwn.net}, very interesting news site
    about Linux in general, and specifically about the kernel. Weekly
    edition, available for free after one week for non-paying
    visitors.
  \end{itemize}


=== International conferences (1)
  
#table(columns: (50\%, 50\%), stroke: none,
[
    \begin{itemize}
    #include "../../common/elc.typ"

    #include "../../common/lpc.typ"

    \end{itemize}
    \column{0.2\textwidth}
    \begin{center}
      \includegraphics[width=\textwidth]{common/elc-logo.png}\\
      \vspace{1cm}
      \includegraphics[width=0.8\textwidth]{common/lpc-logo.jpg}\\
    \end{center}
  \\]


=== International conferences (2)
  
#table(columns: (50\%, 50\%), stroke: none,
[
  \begin{itemize}
  \item FOSDEM: \url{https://fosdem.org}
    \begin{itemize}
    \item Brussels (Belgium), February
    \item Community-oriented conference, free, during the week-end
    \item Many {\em developer rooms}, including on low-level, embedded
      and hardware topics
    \end{itemize}
  \item Embedded Recipes: \url{https://embedded-recipes.org}
    \begin{itemize}
    \item Paris (France), September
    \item 2-day conference about all embedded Linux topics
    \item Well attended by known contributors
    \item Very affordable conference, thanks to sponsors (like Bootlin).
    \end{itemize}
  \item Most conferences are now also accessible on-line, which makes
    them much more affordable.
  \end{itemize}
  \column{0.3\textwidth}
    \begin{center}
      \includegraphics[width=0.8\textwidth]{slides/sysdev-references/fosdem.jpg}\\
      \vspace{1cm}
      \includegraphics[width=0.8\textwidth]{slides/sysdev-references/embedded-recipes.png}\\
    \end{center}
  \\]

